% !TeX program = xelatex
\documentclass[12pt,spanish,letterpaper,oneside]{article}
\PassOptionsToPackage{es-noshorthands}{babel}
\def\documenttitle{Integración de funciones racionales mediante descomposición en fracciones parciales y comprobación diferencial}
\def\documentsubtitle{Evidencia 2 --- Descomposición algebraica y modelos de acumulación (Fase 1)}
\def\documentsubject{Cálculo Integral --- Ingeniero Agrónomo}
\def\documentauthor{Carlos Mauricio Carvajal Coronado}
\def\coursename{Cálculo Integral}
\def\coursecode{REF\#261013755131}
\def\universityname{Universidad Autónoma de Nuevo León}
\def\universityfaculty{Facultad de Agronomía}
\def\universitydepartment{Ingeniero Agrónomo}
\def\universitydepartmentimage{departamentos/logos-institucionales.pdf}
\def\universitydepartmentimagecfg{height=2.1cm}
\def\universitylocation{General Escobedo / Marín, Nuevo León, México}
\def\authortable{%
	\begin{Form}
	\begin{tabular}{@{}ll@{}}
		\textbf{Estudiante:} & Carlos Mauricio Carvajal Coronado\\[0.35cm]
		\textbf{Matrícula:} & 1977875\\[0.35cm]
		\textbf{Docente:} & Cátedra de Matemáticas Aplicadas\\[0.35cm]
		\textbf{Unidad de Aprendizaje:} & Cálculo Integral (Evidencia 2 \textbar\ 15 Puntos)\\[0.5cm]
		\multicolumn{2}{l}{Fecha de realización: 14 de octubre de 2026}\\
		\multicolumn{2}{l}{\universitylocation}
	\end{tabular}
	\end{Form}
}

\usepackage{fontspec}
\setmainfont{Arial}
\setsansfont{Arial}
\AddToHook{file/config.tex/after}[aulatex-calculo-integral]{%
  \def\documentfontsize{12}
  \def\documentinterline{1}
  \def\documentparindent{0}
  \def\documentparskip{5}
  \def\fontdocument{custom}
  \def\fonttypewriter{custom}
  \def\stylecitereferences{custom}
  \def\portraitstyle{style1}
  \def\hfstyle{style1}
  \def\hfwidthwrap{true}
  \def\hfwidthtitle{0.44}
  \def\hfwidthcourse{0.52}
  \def\predocpageromannumber{false}
  \def\pagemarginleft{2.5}
  \def\pagemarginright{2.5}
  \def\pagemargintop{2.5}
  \def\pagemarginbottom{2.5}
  \def\firstpagemargintop{3.8}
  \def\sectionfontsize{\large}
  \def\sectionspacingtop{8}
  \def\sectionspacingbottom{4}
  \def\ssectionfontsize{\normalsize}
  \def\ssectionspacingtop{6}
  \def\ssectionspacingbottom{3}
  \def\defaultimagefolder{UANL/img/}
  \def\usepdfmetadata{false}
}
% Template:     Informe LaTeX
% Documento:    Núcleo del template
% Versión:      DEV
% Codificación: UTF-8
%
% Autor: Pablo Pizarro R.
%        pablo@ppizarror.com
%
% Manual template: [https://latex.ppizarror.com/informe]
% Licencia MIT:    [https://opensource.org/licenses/MIT]

% -----------------------------------------------------------------------------
% CONFIGURACIONES
% -----------------------------------------------------------------------------
% Definiciones previas
\input{src/defs}

% Ajustes usuario
\input{src/config}

% -----------------------------------------------------------------------------
% IMPORTACIÓN DE LIBRERÍAS
% -----------------------------------------------------------------------------
\input{src/env/imports}

% -----------------------------------------------------------------------------
% IMPORTACIÓN DE FUNCIONES Y ENTORNOS
% -----------------------------------------------------------------------------
\input{src/cmd/core}
\input{src/cmd/math}
\input{src/cmd/equation}
\input{src/cmd/image}
\input{src/cmd/title}
\input{src/cmd/other}
\input{src/cmd/column}
\input{src/env/environments}

% -----------------------------------------------------------------------------
% IMPORTACIÓN DE ESTILOS
% -----------------------------------------------------------------------------
\input{src/style/code}
\input{src/style/other}

% -----------------------------------------------------------------------------
% CONFIGURACIÓN INICIAL DEL DOCUMENTO
% -----------------------------------------------------------------------------
\input{src/cfg/init}

% -----------------------------------------------------------------------------
% PORTADA
% -----------------------------------------------------------------------------
\input{src/page/portrait}
\input{src/page/portrait_config}

% -----------------------------------------------------------------------------
% CONFIGURACIÓN DE PÁGINA Y ENCABEZADOS
% -----------------------------------------------------------------------------
\input{src/cfg/page}

% -----------------------------------------------------------------------------
% TABLA DE CONTENIDOS - ÍNDICE
% -----------------------------------------------------------------------------
\input{src/page/index}

% -----------------------------------------------------------------------------
% CONFIGURACIONES FINALES
% -----------------------------------------------------------------------------
\input{src/cfg/final}

\RemoveFromHook{file/config.tex/after}[aulatex-calculo-integral]

\usepackage{tabularx}
\usepackage{booktabs}
\usepackage{amsmath,amssymb}
\usepackage{framed}
\usepackage{tcolorbox}
\usepackage[backend=biber,style=apa,sorting=nyt]{biblatex}
\definecolor{uanlblue}{RGB}{0,55,108}
\definecolor{tableblue}{RGB}{233,240,247}
\def\portraittitlecolor{uanlblue}
\def\sectioncolor{uanlblue}
\def\ssectioncolor{uanlblue}
\def\documenttitlehf{Cálculo Integral: Evidencia 2 --- Fracciones Parciales}
\hypersetup{colorlinks=true,linkcolor=uanlblue,urlcolor=uanlblue,citecolor=uanlblue,pdfauthor={Carlos Mauricio Carvajal Coronado},pdftitle={Integración de funciones racionales mediante fracciones parciales},pdfsubject={Evidencia 2 - Cálculo Integral UANL}}
\addbibresource{UANL/ingeniero-agronomo/calculo-integral/calculo-integral.bib}
\newcolumntype{Y}{>{\raggedright\arraybackslash}X}
\setlength{\emergencystretch}{2.5em}

\let\citep\parencite

\newenvironment{examplebox}[1]{%
  \begin{tcolorbox}[colback=blue!3,colframe=uanlblue,boxrule=0.4pt,
    arc=0pt,left=2mm,right=2mm,top=2mm,bottom=2mm,
    before skip=8pt,after skip=8pt]%
  \setstretch{1.12}%
  \noindent\textbf{#1}\par\smallskip
}{\end{tcolorbox}}

\begin{document}
\hypersetup{pageanchor=false}
\begingroup
\fontsize{14}{17}\selectfont
\renewcommand{\Huge}{\fontsize{14}{17}\selectfont\bfseries}
\renewcommand{\large}{\fontsize{14}{17}\selectfont}
\templatePortrait
\clearpage
\endgroup
\hypersetup{pageanchor=true}
\templatePagecfg
\singlespacing
\setlength{\tabcolsep}{8pt}
\renewcommand{\arraystretch}{1.18}

\begin{abstractd}
Se resuelven las cuatro integrales racionales de la Evidencia 2. En cada ejercicio se factoriza el denominador, se plantea la descomposición en fracciones parciales, se determinan sus coeficientes y se integra término a término. Finalmente, cada antiderivada se comprueba por diferenciación, respetando los valores excluidos del dominio original.
\end{abstractd}

\templateIndex

\section{Introducción}
Una función racional es un cociente de polinomios. Cuando el grado del numerador es menor que el del denominador y este último puede factorizarse, la descomposición en fracciones parciales permite expresar el integrando como una suma de términos elementales. En esta actividad todos los factores son lineales y distintos, por lo que se utiliza el modelo
\[
 \frac{P(x)}{(a x+b)(c x+d)}
 =\frac{A}{a x+b}+\frac{B}{c x+d}.
\]
Después de hallar $A$ y $B$, se aplica
\[
 \int\frac{dx}{a x+b}=\frac1a\ln|a x+b|+C.
\]
La solución es válida en cada intervalo del dominio del integrando original \citep{stewart2016calculus,openstax2016calculus1}.

\section{Procedimiento general}
Antes de aplicar fracciones parciales se comparan los grados del numerador y del denominador. En los cuatro ejercicios el grado del numerador es menor, de modo que las fracciones son propias y no se requiere división polinómica. Para resolver cada integral se siguen cinco pasos:
\begin{enumerate}
\item factorizar completamente el denominador;
\item proponer una fracción simple por cada factor lineal;
\item eliminar denominadores y comparar coeficientes;
\item integrar cada término, incluyendo el factor de la regla de la cadena;
\item derivar el resultado y recuperar el integrando original.
\end{enumerate}

\section{Resolución de ejercicios}

\begin{examplebox}{Ejercicio 7: $\displaystyle \int\frac{1}{x^2-9}\,dx$}
\textbf{1. Factorización y dominio.} Se reconoce una diferencia de cuadrados:
\[
 x^2-9=x^2-3^2=(x-3)(x+3),
 \qquad x\ne-3,\;x\ne3.
\]

\textbf{2. Planteamiento de fracciones parciales.} Como los factores son lineales y distintos,
\[
 \frac{1}{(x-3)(x+3)}=\frac{A}{x-3}+\frac{B}{x+3}.
\]
Al multiplicar por $(x-3)(x+3)$ se obtiene
\[
 1=A(x+3)+B(x-3)=(A+B)x+(3A-3B).
\]
La comparación con $0x+1$ produce el sistema
\[
 \begin{cases}
 A+B=0,\\
 3A-3B=1.
 \end{cases}
\]
De $B=-A$ y $3A-3(-A)=1$ resulta
\[
 6A=1,\qquad A=\frac16,\qquad B=-\frac16.
\]
Por tanto,
\[
 \frac{1}{x^2-9}=\frac1{6(x-3)}-\frac1{6(x+3)}.
\]

\pagebreak
\textbf{3. Integración término a término.} Para la primera integral se toma $u=x-3$, por lo que $du=dx$; para la segunda, $v=x+3$ y $dv=dx$. Así,
\[
 \int\frac{dx}{x-3}=\int\frac{du}{u}
 =\ln|u|+C=\ln|x-3|+C,
\]
\[
 \int\frac{dx}{x+3}=\int\frac{dv}{v}
 =\ln|v|+C=\ln|x+3|+C.
\]
Sustituyendo ambas antiderivadas con sus coeficientes,
\begin{align*}
 \int\frac{1}{x^2-9}\,dx
 &=\frac16\int\frac{dx}{x-3}-\frac16\int\frac{dx}{x+3}\\
 &=\frac16\ln|x-3|-\frac16\ln|x+3|+C\\
 &=\boxed{\frac16\ln\left|\frac{x-3}{x+3}\right|+C}.
\end{align*}
La última igualdad usa $\ln|a|-\ln|b|=\ln|a/b|$.

\textbf{4. Comprobación.}
\begin{align*}
 \frac{d}{dx}\left[\frac16\ln|x-3|-\frac16\ln|x+3|\right]
 &=\frac1{6(x-3)}-\frac1{6(x+3)}\\
 &=\frac{(x+3)-(x-3)}{6(x-3)(x+3)}\\
 &=\frac6{6(x^2-9)}=\frac1{x^2-9}.
\end{align*}
Así se recupera el integrando para $x\ne\pm3$.
\end{examplebox}

\begin{examplebox}{Ejercicio 8: $\displaystyle \int\frac{1}{4x^2-1}\,dx$}
\textbf{1. Factorización y dominio.} El denominador también es una diferencia de cuadrados:
\[
 4x^2-1=(2x)^2-1^2=(2x-1)(2x+1),
 \qquad x\ne-\frac12,\;x\ne\frac12.
\]

\textbf{2. Planteamiento de fracciones parciales.}
\[
 \frac{1}{(2x-1)(2x+1)}
 =\frac{A}{2x-1}+\frac{B}{2x+1}.
\]
Al eliminar denominadores,
\[
 1=A(2x+1)+B(2x-1)=(2A+2B)x+(A-B).
\]
Por comparación de coeficientes,
\[
 \begin{cases}
 2A+2B=0,\\
 A-B=1.
 \end{cases}
\]
La primera ecuación da $B=-A$; al sustituir en la segunda,
\[
 A-(-A)=1,\qquad A=\frac12,\qquad B=-\frac12.
\]
Entonces
\[
 \frac1{4x^2-1}=\frac1{2(2x-1)}-\frac1{2(2x+1)}.
\]

\textbf{3. Integración término a término.} En el primer término se usa
\[
 u=2x-1,\qquad du=2\,dx,
 \qquad dx=\frac12du.
\]
Por ello,
\[
 \int\frac{dx}{2x-1}=\frac12\int\frac{du}{u}
 =\frac12\ln|u|+C=\frac12\ln|2x-1|+C.
\]
En el segundo término, $v=2x+1$, $dv=2\,dx$ y $dx=dv/2$, de modo que
\[
 \int\frac{dx}{2x+1}=\frac12\int\frac{dv}{v}
 =\frac12\ln|2x+1|+C.
\]
Al regresar a la descomposición en fracciones parciales,
\begin{align*}
 \int\frac{1}{4x^2-1}\,dx
 &=\frac12\int\frac{dx}{2x-1}-\frac12\int\frac{dx}{2x+1}\\
 &=\frac12\left(\frac12\ln|2x-1|\right)
   -\frac12\left(\frac12\ln|2x+1|\right)+C\\
 &=\boxed{\frac14\ln\left|\frac{2x-1}{2x+1}\right|+C}.
\end{align*}
Los factores adicionales $1/2$ proceden de los diferenciales, y la última igualdad usa la regla del cociente para logaritmos.

\textbf{4. Comprobación.}
\begin{align*}
 \frac{d}{dx}\left[\frac14\ln|2x-1|-\frac14\ln|2x+1|\right]
 &=\frac1{2(2x-1)}-\frac1{2(2x+1)}\\
 &=\frac{(2x+1)-(2x-1)}{2(2x-1)(2x+1)}\\
 &=\frac2{2(4x^2-1)}=\frac1{4x^2-1}.
\end{align*}
La igualdad vale para $x\ne\pm\frac12$.
\end{examplebox}

\begin{examplebox}{Ejercicio 9: $\displaystyle \int\frac{5}{x^2+3x-4}\,dx$}
\textbf{1. Factorización y dominio.} Se buscan dos números cuyo producto sea $-4$ y cuya suma sea $3$; estos números son $4$ y $-1$. Por ello,
\[
 x^2+3x-4=(x+4)(x-1),
 \qquad x\ne-4,\;x\ne1.
\]

\textbf{2. Planteamiento de fracciones parciales.}
\[
 \frac5{(x+4)(x-1)}=\frac{A}{x+4}+\frac{B}{x-1}.
\]
Se eliminan los denominadores y se agrupan términos:
\[
 5=A(x-1)+B(x+4)=(A+B)x+(-A+4B).
\]
Al comparar con $0x+5$,
\[
 \begin{cases}
 A+B=0,\\
 -A+4B=5.
 \end{cases}
\]
Como $A=-B$, la segunda ecuación queda
\[
 -(-B)+4B=5,\qquad 5B=5,
\]
de donde $B=1$ y $A=-1$. Así,
\[
 \frac5{x^2+3x-4}=-\frac1{x+4}+\frac1{x-1}.
\]

\textbf{3. Integración término a término.} En el primer término se toma $u=x+4$, $du=dx$; en el segundo, $v=x-1$, $dv=dx$. Entonces
\[
 -\int\frac{dx}{x+4}=-\int\frac{du}{u}
 =-\ln|u|+C=-\ln|x+4|+C,
\]
\[
 \int\frac{dx}{x-1}=\int\frac{dv}{v}
 =\ln|v|+C=\ln|x-1|+C.
\]
Al sumar ambas antiderivadas,
\begin{align*}
 \int\frac5{x^2+3x-4}\,dx
 &=-\int\frac{dx}{x+4}+\int\frac{dx}{x-1}\\
 &=-\ln|x+4|+\ln|x-1|+C\\
 &=\boxed{\ln\left|\frac{x-1}{x+4}\right|+C}.
\end{align*}
La forma compacta se obtiene con $\ln|a|-\ln|b|=\ln|a/b|$.

\textbf{4. Comprobación.}
\begin{align*}
 \frac{d}{dx}\left[-\ln|x+4|+\ln|x-1|\right]
 &=-\frac1{x+4}+\frac1{x-1}\\
 &=\frac{-(x-1)+(x+4)}{(x+4)(x-1)}\\
 &=\frac5{x^2+3x-4}.
\end{align*}
La derivada coincide con el integrando para $x\ne-4,1$.
\end{examplebox}

\begin{examplebox}{Ejercicio 10: $\displaystyle \int\frac{x+2}{x^2+11x+18}\,dx$}
\textbf{1. Factorización y dominio.} Se buscan dos números cuyo producto sea $18$ y cuya suma sea $11$; son $2$ y $9$. Entonces
\[
 x^2+11x+18=(x+2)(x+9),
 \qquad x\ne-2,\;x\ne-9.
\]

\textbf{2. Planteamiento de fracciones parciales.}
\[
 \frac{x+2}{(x+2)(x+9)}
 =\frac{A}{x+2}+\frac{B}{x+9}.
\]
Al multiplicar por $(x+2)(x+9)$,
\[
 x+2=A(x+9)+B(x+2)=(A+B)x+(9A+2B).
\]
La comparación de coeficientes conduce a
\[
 \begin{cases}
 A+B=1,\\
 9A+2B=2.
 \end{cases}
\]
De $B=1-A$ y la segunda ecuación se obtiene
\[
 9A+2(1-A)=2,\qquad 7A=0,
\]
por lo que $A=0$ y $B=1$. En consecuencia,
\[
 \frac{x+2}{(x+2)(x+9)}=\frac1{x+9}.
\]
Esta simplificación es válida en el dominio original: aunque el factor $x+2$ se cancela algebraicamente, $x=-2$ continúa excluido.

\textbf{3. Integración.} Se toma $u=x+9$; entonces $du=dx$ y
\[
 \int\frac{dx}{x+9}=\int\frac{du}{u}
 =\ln|u|+C=\ln|x+9|+C.
\]
Por ello,
\begin{align*}
 \int\frac{x+2}{x^2+11x+18}\,dx
 &=\int\frac{dx}{x+9}\\
 &=\boxed{\ln|x+9|+C},
 \qquad x\ne-2,-9.
\end{align*}

\textbf{4. Comprobación.} Se trabaja solamente en el dominio original, es decir, para $x\ne-2,-9$.
\begin{align*}
 \frac{d}{dx}\ln|x+9|
 &=\frac1{x+9}\\
 &=\frac{x+2}{(x+2)(x+9)}\\
 &=\frac{x+2}{x^2+11x+18}.
\end{align*}
La derivada recupera exactamente el integrando en todos los puntos donde este se encuentra definido.
\end{examplebox}

\section{Conclusiones}
Los cuatro ejercicios comparten una estructura racional propia para fracciones parciales, pero cada uno exige cuidar un detalle distinto. Los ejercicios 7 y 8 usan diferencias de cuadrados; el ejercicio 8 añade factores lineales con derivada $2$, lo que introduce coeficientes adicionales al integrar; el ejercicio 9 requiere factorizar un trinomio; y el ejercicio 10 muestra que cancelar un factor no reincorpora al dominio el valor que anulaba el denominador original. En cada caso, la derivación de la respuesta reconstruye exactamente el integrando y confirma tanto los signos como las constantes obtenidas.

\nocite{openstax2016calculus2}


\newpage
\printbibliography[title={Referencias bibliográficas}]
\end{document}

